\begin{abstract}
Feed-forward 3D reconstruction enables direct geometry prediction from single or sparse images, significantly reducing the dependency on per-scene iterative optimization and offering a new paradigm for efficient 3D perception and scene modeling. However, reconstructing transparent objects remains constrained by the inherent ambiguity of single-layer depth representations, where a single pixel often captures both the reflective surface of the glass and the occluded background scene. This ambiguity leads to two fundamental challenges: 1) depth estimation confusion between the transparent surface and the background, and 2) the inability to recover occluded background geometry when refraction occurs. To address these issues, we propose \textbf{Phantom}, a layered geometry reconstruction framework designed for frozen 3D geometry backbones. Phantom employs two complementary image front-ends to decouple the visual signals in the image domain: it independently extracts the interference-free background scene and the opaque surface appearance. We further introduce a geometry-aware loss based on frozen depth networks to suppress subtle artifacts that often induce erroneous surface predictions. Phantom demonstrates exceptional versatility and can be directly integrated with various state-of-the-art geometry backbones (e.g., VGGT, $\pi^3$, DA-V2 Metric, DA3, MoGe-2, and MoGe-3). Experimental results on the ClearGrasp dataset show that Phantom significantly reduces depth errors for transparent surfaces across all backbone models while effectively preserving the underlying background geometry. Finally, to facilitate the learning and evaluation of layered 3D geometry, we introduce, to our knowledge, the first multi-view synthetic dataset that pairs each transparent scene image with camera-aligned ground truth backgrounds, opaque surface references, and dual-layer depth annotations. Our code, pre-trained models, and dataset will be made publicly available.
\end{abstract}
